\section{Jacobi levels and Hill-region topology}
\label{sec:hill}

The Jacobi integral provides the geometric link between the equilibrium
bifurcations of Section~\ref{sec:results} and the periodic dynamics of
Section~\ref{sec:periodic}. We first follow the critical Jacobi values
associated with the equilibrium branches, then resolve one representative
change in Hill-region connectivity in the eleven-equilibrium regime, and
finally place selected periodic families within their corresponding
accessible regions.


\subsection{Critical Jacobi levels along the equilibrium branches}
\label{subsec:critical_jacobi}

Under the convention of Section~\ref{subsec:jacobi}, motion at Jacobi
constant $C_J$ is restricted to the Hill region
\begin{equation}
\mathcal H(C_J)
=
\left\{
(x,y):2\Omega(x,y)\geq C_J
\right\},
\label{eq:hill_region_results}
\end{equation}
whose boundary is the zero-velocity curve
\[
2\Omega(x,y)=C_J.
\]

Away from the primary singularities, the restricted equilibria are the
finite critical points of $\Omega$. Each equilibrium branch $L_i(s)$
therefore determines a critical Jacobi level
\begin{equation}
C_{J,i}(s)
=
2\Omega\!\left(L_i(s)\right).
\label{eq:critical_jacobi_branch}
\end{equation}
For a regular finite level set, changes in the topology of the Hill region
can occur only when $C_J$ passes through such a critical value. Depending
on the local type of the critical point, the change may involve the opening
or closing of a neck, the creation or loss of a bounded component, or a
change in the number of holes; a critical value does not therefore imply a
change in component count in every case.

Figure~\ref{fig:jacobi_levels} shows
$C_{J,i}(s)$ along the reliably tracked equilibrium branches. The
near-primary equilibria generally occupy the upper part of the critical
level ladder, whereas the more distant equilibria occur at lower levels.
Branches created or destroyed at the saddle-node folds introduce,
respectively remove, pairs of critical levels. At the fold itself the two
equilibrium branches coalesce at the same critical point and hence at the
same Jacobi value. The fold cascade of Section~\ref{sec:results} therefore
reorganizes not only the number of equilibria but also the hierarchy of
finite critical energy levels available to the restricted body.

\begin{figure}[htbp]
  \centering
  \includegraphics[width=\textwidth]{critical_jacobi_vs_s.pdf}
  \caption{Critical Jacobi values
  $C_{J,i}(s)=2\Omega(L_i(s))$ along the reliably tracked equilibrium
  branches of the two admissible arcs. Dashed vertical lines mark the
  saddle-node folds. Each curve gives a finite critical level at which the
  regular Hill-region topology may change.}
  \label{fig:jacobi_levels}
\end{figure}


\subsection{A representative connectivity change}
\label{subsec:hill_connectivity}

We resolve a representative connectivity change in the
eleven-equilibrium regime on the Case-2 arc, at
\[
s\simeq0.786,
\qquad
(c_1,c_2)\simeq(0.432,-0.602).
\]
Two Jacobi levels bracketing a portion of the interior critical-level
ladder are shown in the top row of Figure~\ref{fig:hill_regions}. They are
chosen as
\[
C_J=11.63
\qquad\text{and}\qquad
C_J=12.83,
\]
approximately $\pm0.6$ about the median equilibrium critical value
$C_J\simeq12.23$ at this configuration. The interval contains three finite
equilibrium critical values, approximately
\[
11.824,\qquad 11.924,\qquad 12.233.
\]

At $C_J=11.63$ the allowed set forms one connected component on the
computational domain, whereas at $C_J=12.83$ an inner allowed region has
detached from the outer region, giving two connected components. Increasing
$C_J$ across this interval therefore produces a net neck-closing transition
that separates a bounded inner component from the unbounded exterior
component. Because the interval spans a small cluster of critical values,
this calculation demonstrates the net connectivity change across that
portion of the critical-level ladder; it does not assign the transition to
one isolated critical equilibrium.

The component count is evaluated numerically by sampling
$2\Omega(x,y)$ on a uniform
\[
400\times400
\]
grid over
\[
[-6,6]^2,
\]
forming the binary allowed set
\[
2\Omega(x,y)\geq C_J,
\]
and applying four-connectivity labelling. Because
\[
2\Omega(x,y)\rightarrow+\infty
\]
both as a primary is approached and as
$\sqrt{x^2+y^2}\rightarrow\infty$, sufficiently small punctured
neighbourhoods of the primaries and the far exterior belong to the allowed
region for every finite $C_J$. No artificial exclusion mask is introduced
around the primaries in the component calculation.

The calculation was repeated on
\[
300\times300
\qquad\text{and}\qquad
500\times500
\]
grids. The detailed location and shape of the narrow neck are naturally
resolution dependent, but the reported change from one to two connected
components at the two quoted Jacobi levels is unchanged.

This calculation is intentionally representative rather than exhaustive.
Figure~\ref{fig:jacobi_levels} contains many other finite critical levels,
and crossing them need not always alter the component count; a complete
classification would require resolving the topology on both sides of every
critical level. We do not attempt that census here.


\subsection{Periodic families within the accessible regions}
\label{subsec:hill_orbits}

The lower panels of Figure~\ref{fig:hill_regions} place representative
Floquet-stable periodic orbits within the Hill regions corresponding to
their own Jacobi constants. The Case-2 example uses family~3 at
approximately
\[
C_J=21.34,
\]
while the Case-1 example uses family~2 at approximately
\[
C_J=24.04.
\]
In both cases the periodic orbit lies entirely inside a single connected
allowed region and remains separated from the zero-velocity boundary.

For the displayed trajectories we also evaluate the Hill margin
\[
\Delta_H(t)=2\Omega(x(t),y(t))-C_J.
\]
The minimum values along the plotted family-3 and family-2 members are
approximately
\[
0.529
\qquad\text{and}\qquad
2.65\times10^{-3},
\]
respectively. Both remain positive, confirming numerically that the
displayed trajectories remain inside the allowed region and do not
intersect the zero-velocity curve.

The examples illustrate the spatial distinction between the different
dynamical regimes encountered in Section~\ref{sec:periodic}. Periodic
librations associated with the selected outer doubly-elliptic equilibria
occupy broad allowed regions away from the most intricate near-primary
geometry, whereas periodic motion associated with inner
saddle$\times$centre equilibria occurs in a more strongly structured
near-primary region. The Hill geometry does not determine Floquet
stability, but it specifies the energetically accessible part of the plane
within which each periodic family must lie.

Taken together, Figures~\ref{fig:jacobi_levels} and
\ref{fig:hill_regions} illustrate the sequence
\[
\text{equilibrium bifurcation}
\longrightarrow
\text{critical Jacobi level}
\longrightarrow
\text{Hill-region topology}
\longrightarrow
\text{periodic dynamics}.
\]
The saddle-node cascade changes the set of finite critical values
$C_{J,i}(s)$; crossing those values can reorganize the accessible-region
topology; and the periodic families of Section~\ref{sec:periodic} are
confined to the resulting Hill regions. The present computation
demonstrates this connection through a resolution-checked representative
transition rather than a complete topological classification.

\begin{figure}[htbp]
  \centering
  \includegraphics[width=0.95\textwidth]{hill_topology_with_orbits.pdf}
  \caption{Hill regions (shaded: allowed,
  $2\Omega\geq C_J$; blue: zero-velocity curve). Stars denote primaries,
  red points equilibria, and purple curves periodic orbits. Top: the
  eleven-equilibrium regime at $C_J=11.63$ (left) and $C_J=12.83$
  (right), bracketing a small cluster of finite critical Jacobi levels;
  the allowed set changes from one connected component to two across this
  interval as a neck closes.
  Bottom: representative Floquet-stable periodic orbits in the Case-2
  (family~3, left) and Case-1 (family~2, right) regimes, each shown at its
  own Jacobi constant.}
  \label{fig:hill_regions}
\end{figure}